\chapter{AI की नैतिकता एवं भविष्य}
\label{ch:ethics}

\lakshyaheading
\begin{itemize}[leftmargin=1.4em]
  \item AI की सीमाएँ - यह सदैव ``सही'' क्यों नहीं होती।
  \item data में पक्षपात (bias) एवं निजता (privacy) की समस्याएँ।
  \item विज्ञान में AI के उत्तरदायी (responsible) उपयोग के सिद्धांत।
  \item भविष्य के क्षेत्र, करियर (careers) एवं अवसर।
\end{itemize}

\section{AI सदा ``सही'' नहीं होती}

यह समझना अत्यंत आवश्यक है कि AI जादू नहीं है। वह केवल उतना ही ``जानती'' है जितना उसे
data से सिखाया गया है। इसी कारण AI कभी-कभी पूरे आत्मविश्वास के साथ \emph{ग़लत} उत्तर भी
दे देती है। विज्ञान में, जहाँ शुद्धता सर्वोपरि है, इसलिए AI के हर परिणाम की
\term{जाँच एवं पुनरावृत्ति}{verification \& reproducibility} अनिवार्य मानी जाती है।

\begin{itemize}[leftmargin=1.4em]
  \item AI केवल \emph{सहसम्बन्ध (correlation)} देख सकती है; \emph{कारण (causation)}
        सिद्ध करना वैज्ञानिक का कार्य है।
  \item AI का परिणाम तब तक स्वीकार नहीं किया जाता जब तक वह प्रयोग या स्वतंत्र विधि से
        सत्यापित न हो।
\end{itemize}

\begin{jante}
आधुनिक भाषा-आधारित AI कभी-कभी ऐसे ``तथ्य'' गढ़ देती है जो वास्तव में होते ही नहीं -
इसे \emph{hallucination} कहते हैं। इसीलिए किसी भी वैज्ञानिक दावे को मूल स्रोत
(original source) से जाँचना आवश्यक है।
\end{jante}

\section{data में पक्षपात (Bias)}

अध्याय~5 में हमने देखा कि एकपक्षीय data से पक्षपाती परिणाम मिलते हैं। यह केवल तकनीकी
समस्या नहीं, बल्कि एक \emph{नैतिक} समस्या भी है, क्योंकि पक्षपाती AI किसी वर्ग के साथ
अन्याय कर सकती है।

\begin{sochiye}
यदि किसी रोग-निदान AI को केवल एक ही आयु-वर्ग के रोगियों का data दिया जाए, तो वह अन्य
आयु के रोगियों पर अच्छा कार्य क्यों नहीं करेगी? ऐसी स्थिति में किसे हानि हो सकती है,
और इससे कैसे बचा जाए? चर्चा कीजिए।
\end{sochiye}

\section{निजता एवं उत्तरदायित्व (Privacy \& Responsibility)}

विज्ञान, विशेषकर चिकित्सा एवं जीव विज्ञान में, अक्सर व्यक्तिगत data (जैसे चिकित्सा-रिकॉर्ड
या DNA) का उपयोग होता है। ऐसे data की \term{निजता}{privacy} एवं सुरक्षा बनाए रखना
वैज्ञानिक का नैतिक उत्तरदायित्व है। data का उपयोग केवल अनुमति (consent) से एवं निर्धारित
उद्देश्य के लिए ही होना चाहिए।

\subsection{उत्तरदायी AI के सिद्धांत (Responsible AI)}
अनेक देशों एवं संस्थाओं ने अब AI के उपयोग हेतु कुछ साझा सिद्धांत अपनाए हैं:
\begin{itemize}[leftmargin=1.4em]
  \item \term{पारदर्शिता}{Transparency} - AI ने किसी निर्णय तक कैसे पहुँचा, यह
        समझाने योग्य हो।
  \item \term{निष्पक्षता}{Fairness} - किसी समूह के साथ पक्षपात न हो।
  \item \term{जवाबदेही}{Accountability} - परिणाम के लिए उत्तरदायी मनुष्य निर्धारित हो।
  \item \term{निजता}{Privacy} - व्यक्तिगत data सुरक्षित रहे।
\end{itemize}

इन चार स्तम्भों (pillars) पर ही उत्तरदायी AI टिकी होती है (चित्र~\ref{fig:responsible-ai})।

\begin{figure}[h]
\centering
\begin{tikzpicture}[font=\sffamily\small,
    pillar/.style={draw=aiblue, thick, fill=ailite, rounded corners,
      minimum width=2.7cm, minimum height=1.3cm, align=center, inner sep=3pt},
    top/.style={draw=aigreen, very thick, fill=aigreenL, rounded corners,
      minimum width=9.6cm, minimum height=1cm, align=center}]
  \node[top] (roof) {\bfseries उत्तरदायी AI (Responsible AI)};
  \node[pillar, below left=0.9cm and 2.4cm of roof.south] (p1)
      {पारदर्शिता\\\footnotesize Transparency};
  \node[pillar, right=0.3cm of p1] (p2) {निष्पक्षता\\\footnotesize Fairness};
  \node[pillar, right=0.3cm of p2] (p3) {जवाबदेही\\\footnotesize Accountability};
  \node[pillar, right=0.3cm of p3] (p4) {निजता\\\footnotesize Privacy};
  \foreach \p in {p1,p2,p3,p4}
    \draw[aigrey, thick] (\p.north) -- (\p.north |- roof.south);
\end{tikzpicture}
\caption{उत्तरदायी AI के चार स्तम्भ - पारदर्शिता, निष्पक्षता, जवाबदेही एवं निजता।}
\label{fig:responsible-ai}
\end{figure}

\begin{gatividhi}[किसी AI उपकरण की सीमाएँ खोजिए]
\textbf{लक्ष्य:} व्यावहारिक रूप से समझना कि AI कहाँ चूकती है।\\[4pt]
\textbf{चरण:}
\begin{enumerate}[leftmargin=1.6em]
  \item पिछले अध्यायों में बनाए किसी model (जैसे Teachable Machine का पत्ती-वर्गीकरण)
        को जानबूझकर कठिन उदाहरण दिखाएँ - जैसे कम रोशनी, उलटी या आधी ढकी वस्तु।
  \item नोट करें कि किन स्थितियों में model ग़लत या अनिश्चित (uncertain) हुआ।
  \item विचार करें - क्या इसका कारण data की कमी या पक्षपात था?
\end{enumerate}
\textbf{निष्कर्ष:} किसी भी AI को उपयोग करने से पहले उसकी सीमाएँ जानना आवश्यक है।
\end{gatividhi}

\section{भविष्य एवं करियर (Future \& Careers)}

AI एवं विज्ञान के संगम से अनेक नए एवं रोमांचक क्षेत्र बन रहे हैं:
\begin{itemize}[leftmargin=1.4em]
  \item \eterm{Bioinformatics} - जीव विज्ञान एवं data का मेल।
  \item \eterm{Computational Physics} - प्रयोग एवं सिमुलेशन।
  \item \eterm{Data Science} - हर क्षेत्र में data से निर्णय।
  \item चिकित्सा, कृषि (agriculture), मौसम एवं पर्यावरण में AI के अनुप्रयोग।
\end{itemize}

\noindent इन सभी क्षेत्रों का आधार वही विषय हैं जो आप अभी पढ़ रहे हैं - Physics,
Chemistry, Biology एवं Mathematics। इनकी मज़बूत समझ ही भविष्य में AI को सार्थक रूप से
उपयोग करने की कुंजी है।

\begin{jante}
भविष्य की अधिकांश वैज्ञानिक खोजें सम्भवतः मनुष्य एवं AI के \emph{सहयोग} से होंगी -
जहाँ AI विशाल data संभालेगी और मनुष्य कल्पना, प्रश्न एवं नैतिक निर्णय देगा।
\end{jante}

% ---------------- परियोजना ----------------
\begin{pariyojana}[वाद-विवाद - विज्ञान में निर्णय किसका?]
\textbf{उद्देश्य:} AI के नैतिक पहलुओं पर तर्कसंगत ढंग से सोचना एवं चर्चा करना।\\[4pt]
\textbf{विषय:} ``विज्ञान में महत्वपूर्ण निर्णय AI को सौंपे जाने चाहिए या नहीं।''\\[4pt]
\textbf{कार्य (पूरी कक्षा):}
\begin{enumerate}[leftmargin=1.6em]
  \item कक्षा को दो दलों में बाँटें - एक पक्ष में, एक विपक्ष में।
  \item प्रत्येक दल इस पुस्तक के उदाहरणों (AlphaFold, रोग-निदान, bias आदि) का उपयोग
        कर तथ्यों के साथ अपना पक्ष तैयार करे।
  \item बारी-बारी से तर्क प्रस्तुत करें; फिर खुली चर्चा करें।
  \item अंत में मिलकर एक \emph{संतुलित निष्कर्ष} लिखें - AI का उपयोग कब उचित है और
        मनुष्य की भूमिका कहाँ अनिवार्य है।
\end{enumerate}
\textbf{मूल्यांकन:} तर्कों की गुणवत्ता, उदाहरणों के उपयोग एवं निष्कर्ष के संतुलन के
आधार पर।
\end{pariyojana}

% ---------------- सारांश ----------------
\begin{bharatai}
\textbf{NITI Aayog} ने ``Responsible AI for All'' रणनीति प्रस्तुत की है, जो AI के
उपयोग में निष्पक्षता, पारदर्शिता एवं समावेशिता (inclusiveness) पर बल देती है।
\end{bharatai}

\section*{सारांश (Summary)}
\addcontentsline{toc}{section}{सारांश}
\begin{itemize}[leftmargin=1.4em]
  \item AI सदा सही नहीं होती; विज्ञान में उसके हर परिणाम की जाँच एवं सत्यापन आवश्यक है।
  \item AI केवल सहसम्बन्ध (correlation) देखती है; कारण (causation) सिद्ध करना मनुष्य का
        कार्य है।
  \item पक्षपात (bias) एवं निजता (privacy) AI की प्रमुख नैतिक चुनौतियाँ हैं।
  \item उत्तरदायी AI के चार स्तम्भ हैं - पारदर्शिता, निष्पक्षता, जवाबदेही एवं निजता।
  \item भविष्य मनुष्य एवं AI के सहयोग का है, जिसकी नींव Physics, Chemistry, Biology
        एवं Mathematics की समझ है।
\end{itemize}

% ---------------- अभ्यास ----------------
\section*{अभ्यास (Exercises)}
\addcontentsline{toc}{section}{अभ्यास}
\textbf{अ. रिक्त स्थान भरिए:}
\begin{enumerate}[leftmargin=1.6em]
  \item AI द्वारा गढ़े गए झूठे ``तथ्य'' को \rule{2.8cm}{0.4pt} कहते हैं।
  \item AI केवल \rule{2.8cm}{0.4pt} देखती है, जबकि कारण सिद्ध करना मनुष्य का कार्य है।
  \item उत्तरदायी AI का वह सिद्धांत जिसमें निर्णय की व्याख्या सम्भव हो,
        \rule{2.8cm}{0.4pt} कहलाता है।
\end{enumerate}

\textbf{ब. लघु उत्तरीय:}
\begin{enumerate}[leftmargin=1.6em]
  \item ``AI सदा सही नहीं होती'' - दो कारण देकर समझाइए।
  \item उत्तरदायी AI (Responsible AI) के चार सिद्धांत लिखिए।
  \item विज्ञान में डेटा-निजता क्यों महत्वपूर्ण है?
\end{enumerate}

\textbf{स. क्रियात्मक / दीर्घ उत्तरीय:}
\begin{enumerate}[leftmargin=1.6em]
  \item किसी model को कठिन उदाहरण दिखाकर उसकी सीमाएँ खोजने वाली गतिविधि कीजिए और
        अपने अवलोकन लिखिए।
  \item ``विज्ञान में AI - वरदान या चुनौती'' विषय पर एक निबंध लिखिए।
  \item आप भविष्य में AI से जुड़ा कौन-सा वैज्ञानिक क्षेत्र चुनना चाहेंगे, और क्यों?
\end{enumerate}

\begin{mcq}
  \item AI द्वारा आत्मविश्वास से दिया गया ग़लत उत्तर कहलाता है:
    \begin{choices}\item मतिभ्रम (hallucination) \item पक्षपात
      \item समाश्रयण \item प्रशिक्षण\end{choices}
  \item निर्णय की व्याख्या सम्भव होने का सिद्धांत है:
    \begin{choices}\item पारदर्शिता (transparency) \item गोपनीयता
      \item गति \item संपीडन\end{choices}
  \item निम्न में से कौन \emph{उत्तरदायी AI} का सिद्धांत नहीं है?
    \begin{choices}\item निष्पक्षता \item जवाबदेही
      \item डेटा-चोरी \item निजता\end{choices}
  \item विज्ञान में अंतिम निर्णय किसके पास रहना चाहिए?
    \begin{choices}\item AI model \item मनुष्य (वैज्ञानिक)
      \item कंप्यूटर \item इंटरनेट\end{choices}
\end{mcq}

\begin{mukhyashabd}
\textbf{मतिभ्रम (Hallucination)}, \textbf{सहसम्बन्ध बनाम कारण (correlation vs causation)},
\textbf{पक्षपात (bias)}, \textbf{निजता (privacy)}, \textbf{उत्तरदायी AI} -
पारदर्शिता, निष्पक्षता, जवाबदेही, निजता।
\end{mukhyashabd}
\selfcheckqr{https://kkreboot.github.io/ai-in-science-hindi/quiz/ch6.html}

% ============================================================
%  पुस्तक का समापन (closing)
% ============================================================
\section*{आगे की राह (The Road Ahead)}
\addcontentsline{toc}{section}{आगे की राह}

इस पुस्तक में हमने देखा कि कृत्रिम बुद्धिमत्ता कोई अलग ``जादुई'' विषय नहीं, बल्कि
विज्ञान को देखने का एक नया दृष्टिकोण है - जहाँ हर प्रयोग, हर मापन एवं हर खोज के पीछे
\emph{data} है। भौतिकी में pattern, रसायन में अणु, और जीव विज्ञान में DNA - सभी को AI
data के रूप में ``पढ़'' सकती है।

परन्तु सबसे महत्वपूर्ण बात यह है कि AI एक \emph{उपकरण} है, वैज्ञानिक नहीं। जिज्ञासा,
कल्पना, प्रश्न पूछने की कला एवं नैतिक निर्णय - ये सदैव मनुष्य के पास रहेंगे। यदि यह
पुस्तक आपके मन में विज्ञान एवं AI के प्रति एक नई जिज्ञासा जगा सकी, तो इसका उद्देश्य
पूर्ण हुआ।

\begin{center}
\vspace{0.5em}
{\sffamily\bfseries\color{aiblue} सीखते रहिए, प्रयोग करते रहिए, प्रश्न पूछते रहिए।}\\[4pt]
\textit{- पुस्तक समाप्त -}
\end{center}
